\section{Coefficient rings with a one-dimensional tangent space in the special fibre}
\label{sec:curvilinear-special-fibre}

The residue field need not be the prime field for the preceding
criterion to apply.  A Cohen coefficient ring separates the residue
field from the nilpotent direction in the special fibre.  We make
this reduction explicit, including its consequence before completion.
For an Artinian local ring $(C,\mathfrak n,k)$, write
\[
 \operatorname{edim}(C)=\dim_k\mathfrak n/\mathfrak n^2.
\]
Thus $\operatorname{edim}(C)\leq1$ means that its maximal ideal
has at most one generator.  This condition concerns the special
fibre, rather than the embedding dimension of the coefficient ring
itself.

\begin{theorem}\label{thm:curvilinear-special-fibre}
Let $(A,\mathfrak m,k)$ be a complete noetherian local ring of mixed
characteristic $(0,p)$.  Suppose that multiplication by $p$ is
injective on $A$, that $A/pA$ is Artinian, and that
\[
 \operatorname{edim}(A/pA)\leq1.
\]
Then the following conditions are equivalent:
\[
 \Phi_A\text{ is an isomorphism},\qquad
 A\text{ is regular},\qquad A\text{ is a DVR}.
\]
More precisely, there is a complete DVR $O\subset A$, with
uniformizer $p$ and residue field $k$, such that, for
$n=\operatorname{length}_{A/pA}(A/pA)$,
\begin{equation}\label{eq:curvilinear-presentation}
 A\cong O[t]/F,\qquad
 F=t^n+\sum_{i=0}^{n-1}a_it^i,\qquad a_i\in pO.
\end{equation}
If $n=1$, then $A=O$.  If $A$ is singular, then $n\geq2$ and
$a_0\in p^2O$.  Let $r$ be the least $1\leq i<n$ for which
$a_i\notin p^2O$, with $r=\infty$ if there is none, and put
\[
 N=\min(n,pr),\qquad N=n\text{ when }r=\infty.
\]
The degree-$p$ generic specialization kernel contains a subgroup
isomorphic to $\bigoplus_{\mathbb N}\Z/N$ which is split in
$K_0(\Lat_A(S_p))$.  Its retraction is defined on that whole
Grothendieck group and annihilates every Young permutation class.
The same assertion holds for $C_p$.  In the singular case, if $A[1/p]$ is a field,
then $N=n$.
\end{theorem}

The integer $n$ in the presentation is the length of the special
fibre over its own residue field.  When that field is finite of
degree $f$ over $\F_p$, the rank of $A$ over $\Z_p$ is $fn$.
It is the relative degree $n$, rather than $fn$, that appears in
the field-case obstruction.

\subsection{The coefficient presentation}

\begin{lemma}\label{lem:curvilinear-cohen-presentation}
Under the hypotheses of Theorem~\ref{thm:curvilinear-special-fibre},
the presentation \eqref{eq:curvilinear-presentation} exists.
\end{lemma}

\begin{proof}
The Cohen structure theorem supplies a coefficient ring of $A$
\citep[Theorem~10.160.8, Tag~032A]{Stacks2026}.  Since $A$ has
characteristic zero, this coefficient ring is a complete DVR $O$
with uniformizer $p$, and its residue field maps isomorphically
to $k$.  Equivalently, the map from the Cohen DVR into $A$ is
injective: its kernel is an ideal of that DVR, and a nonzero
kernel would contain a power of $p$, contrary to injectivity of
multiplication by $p$ on the nonzero ring $A$.

Put $C=A/pA$.  The map $O/pO\to C$ is injective.  Indeed, a
nonzero residue class from $O/pO$ maps to an element of $A$
whose image in $k$ is nonzero; that element is a unit, and
cannot belong to $pA$.  This identifies $k$ with a coefficient
field of $C$.

The Artinian hypothesis gives an integer $c\geq1$ with
$\mathfrak m^c\subset pA$.  Consequently
\[
 p^jA\subset\mathfrak m^j,\qquad
 \mathfrak m^{cj}\subset p^jA\quad(j\geq1).
\]
The $p$-adic and maximal-ideal-adic topologies on $A$ are
therefore equivalent.  In particular $A$ is complete and
separated for the $p$-adic topology.

Choose a $k$-basis $\bar b_1,\ldots,\bar b_d$ of $C$ and
lifts $b_1,\ldots,b_d\in A$.  These lifts form an $O$-basis
of $A$.  For surjectivity, express any element of $A$ modulo
$pA$ as a $k$-linear combination of the $\bar b_i$, lift its
coefficients to $O$, subtract, and divide the difference by
$p$.  Repeating this operation gives coefficients in $O$ as
convergent $p$-adic series, and represents the original element
as their linear combination of the $b_i$.  For independence,
reduce an equality $\sum_i c_i b_i=0$ modulo $pA$.  The basis
property forces every $c_i\in pO$.  Since $p$ is a
nonzerodivisor on $A$, divide the equality by $p$ and repeat.
Each $c_i$ then belongs to $\bigcap_jp^jO=0$.  Thus $A$ is
finite free over $O$.

Write $\mathfrak n$ for the maximal ideal of $C$.  If
$\operatorname{edim}(C)=0$, Nakayama's lemma gives
$\mathfrak n=0$, so $C=k$ and the preceding basis argument,
with $b_1=1$, gives $A=O$.

Otherwise choose $\bar t\in\mathfrak n$ whose image spans
$\mathfrak n/\mathfrak n^2$.  Nakayama's lemma gives
$\mathfrak n=(\bar t)$.  Let $n$ be the least positive
integer with $\bar t^n=0$.  Repeatedly subtracting a residue
in the coefficient field and dividing by $\bar t$ shows that
\[
 1,\bar t,\ldots,\bar t^{n-1}
\]
span $C$ over $k$.  They are independent: in a nonzero
relation, the lowest nonzero coefficient would express
$\bar t^j$ times a unit as zero for some $j<n$.  This
contradicts the definition of $n$.  Hence
$C\cong k[t]/t^n$, and both its $k$-dimension and its
composition length are $n$.

Choose any lift $t\in\mathfrak m$ of $\bar t$.  The same
basis argument now proves that $1,t,\ldots,t^{n-1}$ is an
$O$-basis of $A$.  There is a unique relation
\[
 t^n+\sum_{i=0}^{n-1}a_it^i=0,\qquad a_i\in O.
\]
Reducing modulo $p$ proves $a_i\in pO$ for every $i$.
The homomorphism $O[T]/(T^n+\sum_i a_iT^i)\to A$ carries
one displayed $O$-basis onto the other and is therefore an
isomorphism.  This proves the lemma.
\end{proof}

\subsection{Regularity and the obstruction}

\begin{proof}[Proof of Theorem~\ref{thm:curvilinear-special-fibre}]
The lemma gives the asserted presentation and proves that $A$
is finite flat over the DVR $O$.  Therefore $\dim A=1$.
When $n=1$ we have $A=O$ and the positive regular-ring theorem
applies.

Suppose $n\geq2$.  The local ring $O[t]/F$ has maximal ideal
$(p,t)$.  The localization of $O[t]$ at $(p,t)$ is a regular
local ring of dimension two.  In its cotangent space, the
image of $F$ is the image of $a_0$: the term $t^n$ and every
$a_it^i$ with $i\geq1$ belong to $(p,t)^2$.  Thus the
quotient has embedding dimension one exactly when
$a_0\notin p^2O$, and embedding dimension two otherwise.
Since it has dimension one, this is exactly the stated
regularity criterion.  A one-dimensional noetherian regular
local ring is a DVR.

In the regular case the positive theorem for regular local
coefficient rings proves that $\Phi_A$ is an isomorphism.
In the singular case,
Theorem~\ref{thm:unramified-distinguished-monic} applies to
the coefficient DVR $O$ and the polynomial $F$.  Its section
and whole-group Young-killing retraction give the claimed
split subgroup, with $N=\min(n,pr)$.  Here $N\geq2$,
because $n\geq2$, $p\geq2$, and every finite $r$ is
positive.  The retraction is nonzero, so $\Phi_A$ is not
surjective.  This proves both implications in the criterion.

It remains to verify the field-case exponent.  Suppose that
$L=A[1/p]$ is a field.  Then $F$ is irreducible over
$K=\operatorname{Frac}(O)$.  Put $s=v_O(a_0)$; the
singular case gives $s\geq2$.  Since $O$ is complete, its
valuation has a unique extension to each finite algebraic
extension of $K$.  One way to see the needed uniqueness is
to take the integral closure in a finite Galois splitting
field.  It is finite by the trace-pairing argument, is local
because $O$ is henselian, and is a DVR.  Its unique valuation
is consequently Galois invariant
\citep[Lemmas~10.153.9 and~10.161.8, Tags~04GM and~032L;
Remark~15.113.6, Tag~09E8]{Stacks2026}.
All $n$ roots of $F$ therefore have the same normalized
valuation $s/n$.  Each product of $n-i$ roots has valuation
$(n-i)s/n$, so the elementary-symmetric expression for $a_i$
gives
\[
 v_O(a_i)\geq \frac{(n-i)s}{n}.
\]
If $v_O(a_i)=1$, then
$i\geq n(1-1/s)\geq n/2$.  Hence $pi\geq n$ for every
such index.  This proves $pr\geq n$, or $r=\infty$, and
therefore $N=n$.  No assumption on the cardinality or
perfection of $k$ was used.
\end{proof}

\subsection{Finite algebras over an arbitrary DVR}

The complete criterion descends along completion for a finite free
algebra over any DVR.  Its hypothesis uses the special fibre modulo
the rational prime, which matters when the given coefficient DVR is
ramified.

\begin{corollary}\label{cor:arithmetic-curvilinear}\label{cor:arbitrary-dvr-curvilinear}
Let $O$ be any mixed-characteristic DVR of residue characteristic
$p$, and let $A$ be a nonzero commutative local algebra finite free
over $O$.  Suppose that
\[
 \operatorname{edim}(A/pA)\leq1.
\]
Then $\Phi_A$ is an isomorphism if and only if $A$ is regular.
If $A$ is singular, its degree-$p$ generic kernel contains a
countable sum $\bigoplus_{\mathbb N}\Z/N$ that is split in the
whole lattice Grothendieck group, with a retraction annihilating
the Young permutation span.  Here $N$ is computed from a Cohen
coefficient presentation of $\widehat A$ as in
Theorem~\ref{thm:curvilinear-special-fibre}.
If $\widehat A[1/p]$ is a field, then
\[
 N=\operatorname{length}_{A/pA}(A/pA).
\]
There is no completeness or residue-field hypothesis on $O$.

In particular, when $v_O(p)=1$, this criterion applies to every
finite free local monogenic $O$-algebra, including arbitrary and
imperfect residue fields.  For a ramified coefficient DVR the
hypothesis remains the displayed condition on $A/pA$.
\end{corollary}

\begin{proof}
Let $\pi$ be a uniformizer of $O$, write $p=u\pi^h$ with $u$
a unit, and let $\mathfrak m$ be the maximal ideal of $A$.
The ring $A$ is noetherian of dimension one.  The quotient
$A/\pi A$ is an Artinian local ring, so
$\mathfrak m^c\subset\pi A$ for some $c$.  It follows that
the $\mathfrak m$-adic and $\pi$-adic topologies on $A$ agree;
they also agree with the $p$-adic topology.  Its completion is
\[
 \widehat A=A\otimes_O\widehat O.
\]
This ring is local and complete, and
$\widehat A/p\widehat A=A/pA$.  In particular the special fibre,
its length, and its embedding dimension are unchanged.
Freeness over $O$ makes $p$ a nonzerodivisor on $A$, and faithful
flatness gives the same conclusion for $\widehat A$.
Both rings have dimension one; their maximal-ideal cotangent
spaces agree.  Thus $A$ is regular if and only if $\widehat A$
is regular, also the general completion criterion
\citep[Lemma~15.44.4, Tag~07NY]{Stacks2026}.

If $A$ is regular, the positive theorem for regular local
coefficient rings gives the assertion.  If it is singular,
Theorem~\ref{thm:curvilinear-special-fibre} gives the stated
split subgroup $H$ of the generic kernel over $\widehat A$.
Let $s:H\to K_0(\Lat_{\widehat A}(S_p))$ and $q$ be its
section and retraction.  Inverting $p$ is the same as inverting
$\pi$, so Theorem~\ref{thm:completion-dvr-kernel} identifies
the completed generic kernel with the one over $A$.
Lift $s$ through that isomorphism to $s_A$, and compose $q$
with completion on the whole group to obtain $q_A$.
Then $q_As_A=1_H$.  Completion preserves every Young lattice,
so $q_A$ still annihilates the Young span.  It is nonzero because
$N\geq2$ in the singular case, proving failure of generation.
The completed-field formula follows from the complete theorem
and the unchanged special-fibre length.

Finally suppose $v_O(p)=1$ and write $k_0=O/pO$.  If $A$ is
monogenic over $O$, then $A/pA$ is a finite-dimensional local
quotient of $k_0[T]$.  Its defining ideal is $(f^e)$ for a
monic irreducible polynomial $f$ and an integer $e\geq1$.
Its maximal ideal is generated by the image of $f$, so its
embedding dimension is at most one.  This argument does not
require $f$ to be separable.  The length over the residue field
$k_0[T]/f$ is $e$, which is the field-case exponent; the absolute
$O$-rank $e\deg f$ need not equal it.  The Cohen coefficient
ring used after completion need not extend the original
coefficient DVR, and no such compatibility is required.
\end{proof}


\begin{remark}
For a ramified coefficient DVR with $v_O(p)\geq2$, the monogenic
algebra $A=O[t]/t^2$ illustrates the distinction between the two
special fibres.  In the regular local ring $O[t]_{(\pi,t)}$, both
$p$ and $t^2$ belong to $(\pi,t)^2$.  Thus $A/pA$ has two
independent cotangent directions, represented by $\pi$ and $t$,
whereas $A/\pi A$ has a principal maximal ideal.  The hypothesis
of Corollary~\ref{cor:arbitrary-dvr-curvilinear} is the former one.
\end{remark}

\begin{example}\label{ex:curvilinear-residue-field}
Let $D=\Z_2[\omega]$ with $\omega^2+\omega+1=0$, and set
\[
 A=D[t]/(t^2+4).
\]
The coefficient DVR $D$ has residue field $\F_4$.  The
ring $A$ has rank four over $\Z_2$, while its special fibre
is $\F_4[t]/t^2$ and has length two over its residue field.
Its maximal ideal is $(2,t)$ and its embedding dimension
is two, so it is singular.  Here $r=\infty$ and $N=2$.
The theorem gives a split countable sum of exact-order-two
classes in the degree-two generic kernel, with a
whole-group Young-killing retraction.  The nontrivial
residue field creates no additional multiplicity in this
cyclic order.
\end{example}
